\documentclass[10pt,a4paper]{amsart}
\usepackage[margin=19mm]{geometry}
\usepackage{amsmath,amssymb,mathtools}
\usepackage[scr=rsfs,cal=euler]{mathalfa}
\usepackage{xfrac}
\usepackage[unicode,hidelinks,bookmarks=false]{hyperref}
\newcommand{\ie}{i.e.\ }
\newcommand{\cf}{cf.\ }
\newcommand{\resp}{respectively\ }
\newcommand{\Subsection}{\subsection}
\providecommand{\nobreakdash}{\nobreak\hbox{-}}
\setlength{\emergencystretch}{2em}
\multlinegap0pt
\usepackage{mathtools}
\usepackage[shortlabels]{enumitem}
\usepackage{amssymb}
\usepackage{xcolor}
\newtheorem{lem}{Lemms}
\newcommand{\R}{{\mathbb R}}
\def\bp{\boldsymbol{p}}
\def\bx{\boldsymbol{x}}
\def\by{\boldsymbol{y}}
\def\hB{{\widehat{B}}}
\def\hT{{\widehat{T}}}
\def\one{\mathbf{1}}
\title{Minimal Fillings of Finite Metric Spaces and Convex Polyhedra}
\author{A.O. Ivanov, A.A. Tuzhilin}
\date{Source: arXiv:2607.27211v1 (CC BY 4.0)}
\begin{document}
\maketitle

\noindent\emph{Source and licence.} Minimal Fillings of Finite Metric Spaces and Convex Polyhedra. Version arXiv:2607.27211v1 (CC BY 4.0), distributed by arXiv under the Creative Commons Attribution 4.0 International licence, http://creativecommons.org/licenses/by/4.0/, as recorded in the arXiv licence metadata for this exact version and shown on the abstract page https://arxiv.org/abs/2607.27211v1. Full licence notice: \texttt{licenses/arxiv-2607.27211-CC-BY-4.0.txt}.\par\smallskip

\noindent\textit{Editorial scope.} Counted unit: the article's lem lem:clp\_glp (source0.tex lines 278--280). The article's complete original proof of it is delivered with it (source0.tex lines 282--293, one \texttt{proof} environment of 12 lines). No mathematics is written, repaired or re-derived: the statement and the proof are the article's own lines, in the article's own notation, and no retained line is substituted or re-flowed.  No further source-local context is needed: every cross-reference of every retained line resolves inside the two delivered spans.  Nothing was omitted from any retained span, so the package carries no omission record.  No retained line contains a citation, so the wrapper carries no bibliography and no bibliography entry.  No retained line contains a figure environment, an image-inclusion command or drawing code, so no image is delivered and the package declares no asset dependency.  Editorial formatting only, in addition: an AMS \texttt{amsart} preamble that restates the article's own declarations and macros used by the retained lines, namely the environments \texttt{lem} on separate counters, together with the standard packages those lines need.  The retained environments print in this document as Lemma 1 in order of appearance, whereas the article prints the same items with the numbering of its own full document.  The complete native source of the article as distributed in the arXiv e-print for this exact version is bundled unchanged as \texttt{sources/206335/source0.tex}.
\begin{lem}\label{lem:clp_glp}
If $\bx\in\R^N$ belongs to the polyhedron $T_n$ defined by the system of inequalities $B\bx\ge\one_C$, $\bx\ge0$, then $(\bx,\by)$, where $\by=B\bx -\one_C$, belongs to $\hT_n$, that is, $\hB(\bx,\by)=\one_C$, $\bx\ge0$ and $\by\ge0$. Conversely, if $(\bx,\by)\in\R^{N+C}$ belongs to the polyhedron $\hT_n$ defined by the system of equations and inequalities $\hB(\bx,\by)=\one_C$, $\bx\ge0$ and $\by\ge0$, then $\by=B\bx -\one_C$ and $\bx\in T_n$, that is $B\bx\ge\one_C$, $\bx\ge0$. The projection map $\pi\:\hT_n\to T_n$, $(\bx,\by)\mapsto\bx$ is one-to-one, and the inverse map is $\bx\mapsto(\bx,B\bx-\one_C)$. Moreover, $\bx^*$ is a corner point of $T_n$ if and only if $(\bx^*,\by^*)$, $\by^*=B\bx^*-\one_C$, is a corner point of $\hT_n$.
\end{lem}

\begin{proof}
If $\bx\in T_n$, then $\bx\ge0$ and $\by=B\bx -\one_C\ge0$ by definition of the polyhedron $T_n$, therefore $\hB(\bx,\by)=B\bx-\by=\one_C$, whence $(\bx,\by)\in\hT_n$. Conversely, if $(\bx,\by)\in\hT_n$, then $\bx\ge0$, $\by\ge0$ and $\hB(\bx,\by)=B\bx-\by=\one_C$, whence $B\bx-\one_C=\by\ge0$, that is, $B\bx\ge\one_C$, and, therefore, $\bx\in T_n$.

Further, from the already proven first two statements of Lemma it follows that for any point $(\bx,\by)\in\hT_n$ the inclusion $\bx\in T_n$ holds, and conversely, $(\bx,B\bx-\one_C)\in\hT_n$ holds for each $\bx\in T_n$, and for any point $(\bx,\by)\in\hT_n$ the equality $\by=B\bx-\one_C$ holds, that is, $\by$ is uniquely determined by $\bx$, which means that the projection $\pi$ is bijective.

Now let $\bx^*$ be a corner point of $T_n$, and let $(\bx^*,\by^*)$, where $\by^*=B\bx^*-\one_C$, be not a corner point for $\hT_n$. The latter means that in $\hT_n$ there exist two distinct points $\bp_0=(\bx_0,\by_0)$ and $\bp_1=(\bx_1,\by_1)$ such that $(\bx^*,\by^*)=(1-t)\bp_0+t\bp_1$ for some $t\in(0,1)$. But then $\bx^*=(1-t)\bx_0+t\bx_1$ for the same value of $t$, and $\bx_0$ and $\bx_1$, as we have already proved, belong to $T_n$, therefore $\bx^*$ is not a corner point of $T_n$, a contradiction. Conversely, let $(\bx^*,\by^*)$ be a corner point of $\hT_n$, and let $\bx^*$ not be a corner point of $T_n$. Then there exist points $\bx_0$ and $\bx_1$ in $T_n$ for which $\bx^*=(1-t)\bx_0+t\bx_1$ for some $t\in(0,1)$. Let's put $\by_0=B\bx_0-\one_C$ and $\by_1=B\bx_1-\one_C$. Then $\by_0\ge0$ and $\hB(\bx_0,\by_0)=\one_C$, so $\bp_0=(\bx_0,\by_0)\in\hT_n$, and, in the same way, $\bp_1=(\bx_1,\by_1)\in\hT_n$. Besides,
\begin{multline*}
(1-t)\by_0+t\by_1=(1-t)(B\bx_0-\one_C)+t(B\bx_1-\one_C)=\\
=B\big((1-t)\bx_0+t\bx_1\big)-\one_C=B\bx^*-\one_C=\by^*,
\end{multline*}
from where $(\bx^*,\by^*)=(1-t)\bp_0+t\bp_1$ for the same value of $t\in(0,1)$, that is, $(\bx^*,\by^*)$ is not a corner point for $\hT_n$, a contradiction.
\end{proof}

\end{document}
